% Generated draft: requires mathematical and source review.
\documentclass[11pt]{article}
\usepackage[margin=25mm]{geometry}
\usepackage{fontspec}
\setmainfont{Latin Modern Roman}
\usepackage{amsmath,amssymb,mathtools,unicode-math}
\setmathfont{Latin Modern Math}
\usepackage{enumitem,needspace,xurl,hyperref}
\setlength{\parindent}{0pt}
\setlength{\parskip}{5pt}
\setlength{\emergencystretch}{3em}
\newcommand{\N}{\mathbb{N}}
\newcommand{\Z}{\mathbb{Z}}
\newcommand{\Q}{\mathbb{Q}}
\newcommand{\R}{\mathbb{R}}
\newcommand{\C}{\mathbb{C}}
\newcommand{\F}{\mathbb{F}}
\newcommand{\A}{\mathbb{A}}
\newcommand{\PP}{\mathbb{P}}
\newcommand{\E}{\mathbb{E}}
\newcommand{\eps}{\varepsilon}
\newcommand{\dd}{\,\mathrm{d}}
\begin{document}
\section*{UnsolvedMath ID: 30000512}
\textbf{Source problem code:} catalogue entry
\subsection{Definitions and mathematical statement}
The question compares a threshold at the origin in characteristic zero with the threshold of reductions of the same polynomial. An explicit standard arithmetic setting is a nonzero polynomial $f\in\Z[1/N][x_1,\ldots,x_n]$ with $f(0)=0$, viewed over $\C$. For primes $p\nmid N$ where its reduction is nonzero, let $f_p\in\F_p[x_1,\ldots,x_n]$, let $R_p=\F_p[x_1,\ldots,x_n]_{\mathfrak m}$, and let $\mathfrak m=(x_1,\ldots,x_n)R_p$.

Write $\operatorname{lc}_0(f)$ for the log canonical threshold at $0$. In the complex analytic convention it is
\[
\operatorname{lc}_0(f)=\sup\{c\geq0:|f|^{-2c}\text{ is integrable on some neighborhood of }0\}.
\]
For $q=p^e$, the Frobenius power is $\mathfrak m^{[q]}=(x_1^q,\ldots,x_n^q)R_p$. Set
\[
\nu_{f_p}^{\mathfrak m}(q)=\max\{r\in\Z_{\geq0}: f_p^r\notin\mathfrak m^{[q]}\},\qquad
c^{\mathfrak m}(f_p)=\lim_{e\to\infty}\frac{\nu_{f_p}^{\mathfrak m}(p^e)}{p^e}.
\]
This is the usual maximal-ideal $F$-threshold in this regular local setting.

The requested assertion is that, for each fixed polynomial in the intended arithmetic class,
\[
\#\{p\text{ of good reduction}:\operatorname{lc}_0(f)=c^{\mathfrak m}(f_p)\}=\infty.
\]
The polynomial is fixed before the infinitely many primes are chosen. Convergence of the characteristic-$p$ thresholds to the characteristic-zero threshold does not itself give equality for infinitely many primes.

The supplied extract does not specify the coefficient ring, an integral model, nonzero and origin-vanishing assumptions, or the ideal $\mathfrak m$. The displayed definitions make the customary origin/maximal-ideal interpretation explicit; they do not assert that these missing hypotheses were verified in the supplied record. Polynomials over more general characteristic-zero fields require an identified model and a precise choice of reduction places.
\subsection{Short English statement}
For each fixed polynomial, ask whether its log canonical threshold at the origin agrees exactly with the maximal-ideal $F$-threshold of its reduction for infinitely many primes. Exact equality, rather than convergence, is the target; the arithmetic model and local conventions are missing from the extract.
\subsection{Sources}
\begin{itemize}
\item Supplied report: \emph{Classical Algebraic Geometry} (2007), Oberwolfach Reports, DOI 10.4171/owr/2006/27, pp. 1615--1662; proposed year 2006. \url{https://doi.org/10.4171/owr/2006/27}.
\item Supplied reference labelled ``$F$-threshold limit-point results'': \url{https://arxiv.org/abs/1106.0207}.
\item Local UnsolvedMath registry, permanent ID 30000512, catalogue entry; created 25 August 2026, updated 16 September 2026. It reports a corrected\_verified comparison with the official report. The supplied 21 August 2026 literature triage labels the assertion open, distinguishing known discreteness, rationality and convergence results from the infinitely-many-primes equality question. This is a dated source assessment, not an independent status determination.
\end{itemize}
\end{document}
